% ID: SEL-D0039-OA-E03
% ESTADO: cotejada. Segunda pasada visual y borrador compilado revisados.
% Fuente: Selectividad/matematicas_aplicadas/Ua4MATEMÁTICAS APLICADAS examen.pdf
% SHA-256: 78ff6bdd48cbac71e6cc47d6ddf3dcb9492540ecb9cbe2032ebfec02e39fb66d
% Página: 1; opción A; ejercicio original 3.

De los 700 alumnos matriculados en una asignatura, 210 son hombres y 490 mujeres. Se sabe que el 60\% de los hombres y el 70\% de las mujeres aprueban dicha asignatura. Se elige una persona al azar.
\begin{enumerate}[label=\alph*)]
\item \textbf{(1.5 puntos)} \textquestiondown{}Cuál es la probabilidad de que apruebe la asignatura?
\item \textbf{(1 punto)} Sabiendo que ha aprobado la asignatura, \textquestiondown{}cuál es la probabilidad de que sea una mujer?
\end{enumerate}
